\chapter{Brucellosis}

\section*{Definition and Overview}
Brucellosis is a bacterial infection caused by organisms of the genus \textit{Brucella}, transmitted from infected animals or their products to humans. It is also called undulant fever because of its characteristic relapsing pattern of fevers. People usually acquire the infection by drinking unpasteurised milk or eating dairy products from infected cattle, goats, sheep, or pigs, or through direct contact with infected animals, their blood, or tissues. Brucellosis typically produces a prolonged febrile illness with joint pain, sweating, and fatigue, and it can involve many organs if untreated. It is a notifiable disease in many countries because of its importance in livestock and its public health implications.

\section*{Epidemiology}
Brucellosis occurs worldwide but is most common in the Mediterranean region, the Middle East, parts of Central and South America, Africa, and Central Asia, where brucellosis in livestock is not well controlled. It is primarily an occupational disease of farmers, shepherds, abattoir workers, veterinarians, and laboratory staff, but travellers and consumers can become infected through unpasteurised dairy products. Men are affected more often than women in many regions, largely reflecting greater occupational contact with livestock. The reported incidence falls sharply in countries that implement animal vaccination and milk pasteurisation programmes, demonstrating that the disease is highly preventable.

\section*{Aetiology and Pathophysiology}
The infection is caused by several species of \textit{Brucella}, small Gram-negative bacteria that survive and multiply inside the body's white blood cells, allowing them to evade immune clearance.

\begin{itemize}
    \item \textbf{\textit{B.~melitensis:}} the most pathogenic species, transmitted from goats and sheep, and the commonest cause of severe human disease
    \item \textbf{\textit{B.~abortus:}} transmitted from cattle and associated with milder disease
    \item \textbf{\textit{B.~suis:}} transmitted from pigs and other animals
    \item \textbf{\textit{B.~canis:}} transmitted from dogs, causing illness that is frequently mild or subclinical
\end{itemize}

The bacteria enter the body through the digestive tract, through breaks in the skin, or via the respiratory tract, and are then carried in the bloodstream to organs rich in immune tissue, including the liver, spleen, lymph nodes, and bone marrow. There they multiply within cells, causing granulomatous inflammation and the systemic features of the illness, and their intracellular persistence underlies the chronic, relapsing course of untreated disease.

\section*{Clinical Features}
Symptoms usually begin one to four weeks after exposure and develop gradually over days to weeks.

\begin{itemize}
    \item Fever, often with profuse sweats and chills that come and go (undulant fever)
    \item Malaise, weakness, and profound fatigue
    \item Headache and generalised muscle aches
    \item Joint and back pain, particularly in the lower back, hips, and knees
    \item Weight loss and loss of appetite
    \item Enlargement of the liver or spleen on examination
    \item In some cases, a dry cough, abdominal pain, or testicular swelling
\end{itemize}

\section*{Differential Diagnosis}
Fever, fatigue, and joint pain are features of many conditions, so brucellosis is easily missed without a careful exposure history.

\begin{itemize}
    \item \textbf{Tropical infections:} malaria, typhoid fever, and other causes of prolonged fever
    \item \textbf{Viral illness:} influenza and other systemic viral infections
    \item \textbf{Tuberculosis:} chronic infection with fever, sweats, and weight loss
    \item \textbf{Infective endocarditis:} fever with heart murmur and systemic emboli
    \item \textbf{Connective tissue disease:} systemic lupus erythematosus and other rheumatological conditions
    \item \textbf{Malignancy:} lymphoma and other cancers presenting with fever and night sweats
    \item \textbf{Joint inflammation:} rheumatoid arthritis and other causes of inflammatory arthritis
\end{itemize}

\section*{When to Seek Medical Care}
See a doctor if you develop fever, sweats, joint pain, or fatigue and may have been exposed to animals or unpasteurised dairy products, or if you have travelled to a region where brucellosis is common. Mention this exposure to the doctor, because it is the key to making the diagnosis.

\textbf{Call 000 (or local emergency services) for:}
\begin{itemize}
    \item High fever with confusion or severe drowsiness
    \item Severe headache with neck stiffness
    \item Chest pain, palpitations, or shortness of breath
    \item Sudden severe abdominal pain
    \item Any sign that the infection may have reached the brain, heart, or another vital organ
\end{itemize}

\section*{Diagnosis and Investigations}
Brucellosis is confirmed by tests that detect the bacteria directly or the immune response to them.

\begin{itemize}
    \item \textbf{History:} careful questioning about animal contact, unpasteurised food or drink, occupation, and travel
    \item \textbf{Blood tests:} full blood count, liver function tests, and inflammatory markers
    \item \textbf{Blood cultures:} taken repeatedly to grow the organism, although it is slow-growing and may take several weeks to appear
    \item \textbf{Serology:} antibody testing, which is the mainstay of diagnosis
    \item \textbf{Polymerase chain reaction (PCR):} rapid detection of the organism's genetic material
    \item \textbf{Imaging:} X-rays, ultrasound, or MRI of joints, the spine, or the abdomen when complications are suspected
\end{itemize}

\section*{Management}
Brucellosis requires prolonged treatment, and specialist infectious-disease advice is usually sought.

\begin{itemize}
    \item \textbf{Combination antibiotics:} two agents given together for six weeks or longer, because the intracellular location of the bacteria means single-drug regimens frequently fail
    \item \textbf{Tailored therapy:} the choice and duration of antibiotics depend on the severity and site of infection and on the presence of complications such as endocarditis or neurobrucellosis
    \item \textbf{Symptom relief:} paracetamol or anti-inflammatory medicines for fever and pain
    \item \textbf{Rest and nutrition:} adequate rest and a balanced diet while the illness resolves
    \item \textbf{Monitoring:} follow-up blood tests and clinical review to detect relapse
    \item \textbf{Treatment of complications:} drainage of abscesses or surgery for infected heart valves and spinal collections where required
\end{itemize}

\section*{Complications}
Without correct, prolonged treatment, brucellosis can localise to particular organs and become chronic.

\begin{itemize}
    \item Arthritis, particularly of the sacroiliac joints, hips, and knees
    \item Spondylitis and spinal epidural abscesses
    \item Epididymo-orchitis (infection of the testicle and epididymis)
    \item Endocarditis, the most dangerous complication, which may require surgery
    \item Meningitis and other forms of neurobrucellosis
    \item Hepatitis and liver abscesses
    \item Chronic fatigue, arthralgia, and depression lasting many months
    \item Relapse of symptoms, sometimes after apparently successful treatment
\end{itemize}

\section*{Prognosis and Outcomes}
With early diagnosis and a full course of combination antibiotics, most people recover completely. Symptoms such as fatigue and joint pain may persist for months, and relapse occurs in roughly 5--10\% of treated patients, usually when the course of treatment was too short. Complications such as endocarditis and neurobrucellosis carry a more guarded outlook. In otherwise healthy people who receive proper treatment, brucellosis is rarely fatal.

\section*{Prevention and Self-Care}
Brucellosis is preventable by avoiding exposure to the bacteria.

\begin{itemize}
    \item Drink only pasteurised milk and dairy products, especially when travelling
    \item Cook meat thoroughly and avoid raw or undercooked animal products
    \item Wear gloves, goggles, and protective clothing when handling animals, placentas, or carcasses
    \item Cover cuts and abrasions when working with livestock or animal products
    \item Support control of brucellosis in livestock through vaccination and testing of herds
    \item Follow strict laboratory safety protocols if working with \textit{Brucella} cultures
    \item Maintain high hygiene standards on farms and in slaughterhouses, and keep raw meat separate from other foods
\end{itemize}

\section*{Sources and Further Reading}
The following organisations provide reliable information on brucellosis for patients, travellers, and clinicians.

\begin{itemize}
    \item Centers for Disease Control and Prevention. \textit{Brucellosis: signs, symptoms, and treatment information.} https://www.cdc.gov/brucellosis/
    \item World Health Organization. \textit{Brucellosis: fact sheets and public health guidance.} https://www.who.int/
    \item NHS. \textit{Brucellosis: causes, symptoms, and treatment.} https://www.nhs.uk/conditions/brucellosis/
    \item MedlinePlus. \textit{Brucellosis health information.} https://medlineplus.gov/brucellosis.html
    \item MSD Manual. \textit{Brucellosis: professional overview and management.} https://www.msdmanuals.com/
    \item Department of Health and Aged Care. \textit{Notifiable disease information for health professionals.} https://www.health.gov.au/
\end{itemize}
